\documentclass[uplatex,dvipdfmx,a4paper,11pt,nomag]{jsarticle}
\usepackage[margin=22mm,top=24mm,bottom=24mm]{geometry}
\usepackage[dvipdfmx]{graphicx}
\usepackage[dvipdfmx]{xcolor}
\usepackage{array,booktabs,longtable,tabularx,multirow}
\usepackage{enumitem}
\usepackage[dvipdfmx,unicode,hidelinks]{hyperref}
\usepackage{pxjahyper}
\usepackage{fancyhdr}
\usepackage{titlesec}
\usepackage{titletoc}
\usepackage{framed}
\usepackage{lastpage}
\usepackage{setspace}
\usepackage{ragged2e}
\setstretch{1.18}
\emergencystretch=3em
\sloppy

\definecolor{shadecolor}{gray}{0.96}

\pagestyle{fancy}
\fancyhf{}
\fancyhead[L]{\small 第 05 章 ソフトウェアテスト}
\fancyhead[R]{\small SWEBOK v4.0a 日本語まとめ}
\fancyfoot[C]{\small \thepage/\pageref{LastPage}}
\renewcommand{\headrulewidth}{0.4pt}

\titleformat{\section}{\Large\bfseries}{\thesection}{1em}{}
\titleformat{\subsection}{\large\bfseries}{\thesubsection}{1em}{}
\titleformat{\subsubsection}{\normalsize\bfseries}{\thesubsubsection}{1em}{}
\setcounter{tocdepth}{2}
\setcounter{secnumdepth}{3}
\renewcommand{\contentsname}{目次}
\renewcommand{\figurename}{図}
\renewcommand{\tablename}{表}

\newcolumntype{Y}{>{\raggedright\arraybackslash}X}
\newcolumntype{P}[1]{>{\raggedright\arraybackslash}p{#1}}

\newenvironment{pointbox}[1]{\par\medskip\begin{shaded}\noindent\textbf{#1}\par\smallskip\noindent}{\end{shaded}\par\medskip}
\newenvironment{todobox}[2]{\par\medskip\begin{framed}\noindent\textbf{TODO（図） #1}\par\smallskip\noindent\textbf{画像生成用プロンプト}\par\smallskip\noindent #2\par\smallskip}{\end{framed}\par\medskip}
\newcommand{\term}[2]{\textbf{#1}（#2）}

\begin{document}

\begin{titlepage}
\centering
\vspace*{18mm}
{\LARGE\bfseries 第 05 章 ソフトウェアテスト\par}
\vspace{6mm}
{\large Software Testing の日本語まとめ\par}
\vspace{5mm}
{\normalsize SWEBOK Guide v4.0a Chapter 05 を初学者向けに再構成\par}
\vspace{7mm}
{\normalsize 作成日：2026 年 4 月 29 日\par}
\vspace{10mm}
\hrule
\vspace{3mm}
\begin{flushleft}
\textbf{この資料の主張}\par
ソフトウェアテストは、完成後に不具合を探す作業だけではない。テストは、有限個のテストケースを選び、テスト対象システムを実行し、期待される振る舞いと実際の振る舞いを比べることで、品質に関する根拠を得る活動である。どのテストケースを選ぶか、どの水準で十分とみなすか、結果をどう判断するかによって、テストの有効性は大きく変わる。したがって、テストは技法、測定、プロセス、開発ライフサイクル、ツールを組み合わせて設計する必要がある。\par
\vspace{3mm}
\begin{pointbox}{この資料の読み方}
専門用語は原則として日本語で示し、必要に応じて英語を括弧内に併記する。英語名は、元資料、試験範囲、ツール名を確認するときの手がかりとして残す。初学者が前提知識なしに読めるように、各節では「何を確認する活動か」「なぜその考え方が必要か」「実務では何に注意するか」を重視する。文体は、だである調に統一する。
\end{pointbox}
\end{flushleft}
\vfill
\end{titlepage}

\tableofcontents
\clearpage

\section*{全体概要：この章で学ぶこと}
\addcontentsline{toc}{section}{全体概要：この章で学ぶこと}
この章は、ソフトウェアテストを「実行による確認」として扱う。テストは、テスト対象システムを動かし、観測された結果と期待される結果を比べる活動である。ただし、現実のソフトウェアでは、すべての入力、すべての状態、すべての環境を試すことはできない。したがって、テストは常に、限られた資源でどのテストケースを選ぶかという判断を伴う。

\begin{pointbox}{到達目標}
この章を読み終えると、次のことができるようになる。
\begin{itemize}[leftmargin=2em]
  \item 欠陥、障害、テストケース、テストスイート、テストオラクルを区別できる。
  \item 単体テスト、結合テスト、システムテスト、受入テストの違いを説明できる。
  \item 仕様ベース、構造ベース、経験ベース、欠陥ベース、利用ベースのテスト技法を比較できる。
  \item カバレッジ、欠陥密度、信頼性成長モデル、ミューテーションスコアなどの測度の役割を説明できる。
  \item テスト計画、設計、環境構築、実行、インシデント報告、完了判断の流れを説明できる。
  \item 従来型開発、アジャイル、DevOps、TDD、クラウド、AI、ブロックチェーンなどでテストがどう変わるかを説明できる。
\end{itemize}
\end{pointbox}

初学者は、まず「テストは何を証明でき、何を証明できないか」を押さえるとよい。次に、テストのレベル、目的、技法を分けて理解する。最後に、測定、プロセス、ツールの話に進むと、個別技法が実務の中でどう使われるかを理解しやすい。

\begin{todobox}{05 章の全体地図を入れる}{白背景、モノクロ、A4 本文幅に収まる横長の学習用図解を作成する。中央に「ソフトウェアテスト」を置き、周囲に「基礎」「テストレベル」「テスト技法」「テスト関連測度」「テストプロセス」「開発プロセスと適用領域」「新興技術」「テストツール」を配置する。矢印は「目的を決める」「テストを選ぶ」「実行する」「測る」「改善する」の流れで表す。ラベルは日本語、細線、講義ノート風、余白を広めにする。}
\end{todobox}

\section*{最初に必要な用語}
\addcontentsline{toc}{section}{最初に必要な用語}
\begin{longtable}{P{0.23\linewidth}P{0.69\linewidth}}
\toprule
\textbf{用語} & \textbf{初学者向けの説明} \\
\midrule
テスト対象システム SUT & テストされる対象である。プログラム、アプリケーション、サービス、マイクロサービス、ミドルウェア、システム、システム群、エコシステムまで含み得る。\\
テストケース & 入力値だけでなく、実行条件、時刻条件、手順、期待結果を含む、テスト実行のための具体的な指定である。\\
テストスイート & 複数のテストケースをまとめた集合である。\\
期待結果 & テストを実行したときに、本来得られるべき出力、状態変化、メッセージ、振る舞いである。\\
テストオラクル & 実際の結果が正しいかどうかを判定する基準または仕組みである。人、仕様、モデル、コード注釈、比較ツールなどがなり得る。\\
欠陥 & 障害を引き起こす原因である。コード、設計、要求、設定、運用手順に存在し得る。\\
障害 & 利用者や外部から観測される望ましくない振る舞いである。欠陥が存在しても、条件がそろわなければ障害として現れないことがある。\\
カバレッジ & テストが対象のどの部分をどれだけ通ったかを表す度合いである。文、分岐、条件、要求、状態など、対象は目的によって変わる。\\
回帰テスト & 変更後に、以前できていたことが壊れていないかを再確認するテストである。\\
シフトレフト & テストや品質確認を開発の後半に置くだけでなく、早い段階から行う考え方である。\\
\bottomrule
\end{longtable}

\section*{導入：ソフトウェアテストとは何か}
\addcontentsline{toc}{section}{導入：ソフトウェアテストとは何か}
ソフトウェアテストとは、通常は無限に近い実行領域から有限個のテストケースを選び、テスト対象システムが期待される振る舞いを提供するかを動的に確認する活動である。この定義には、いくつかの重要な意味が含まれる。

第一に、テストは動的である。つまり、対象を実行して確認する。静的解析、レビュー、形式的証明などは品質向上に重要であるが、この章では主に実行を伴うテストを扱う。

第二に、テストは有限である。すべての入力値、すべての状態、すべての環境を試すことは現実的でない。したがって、テストでは、限られた時間、費用、人員の中で、どのテストケースを選ぶかを決める必要がある。

第三に、テストには期待結果が必要である。実行結果を見ても、それが正しいかどうかを判断できなければテストにならない。期待結果を判断する仕組みがテストオラクルである。

\begin{todobox}{テストの基本定義図を入れる}{白背景、モノクロ、A4 本文幅に収まる横長の概念図を作成する。左に「テストケース集合」、中央に「テスト対象システム SUT」、右に「実際の結果」、その下に「期待結果」と「テストオラクル」を配置する。テストケースが SUT に入力され、実際の結果と期待結果をオラクルが比較し、「合格」「不合格」「判定不能」を出す流れを矢印で示す。日本語ラベル、細線、講義ノート風。}
\end{todobox}

\section{ソフトウェアテストの基礎}
\subsection{欠陥と障害}
欠陥は、望ましくない振る舞いの原因である。障害は、その原因が実行時に表に出て、利用者や外部システムから観測される現象である。欠陥があっても、該当する入力や状態に到達しなければ障害として現れないことがある。逆に、障害が観測されても、どの欠陥が原因かを一意に決められるとは限らない。

テストが直接明らかにするのは、多くの場合、障害である。テストに失敗した後は、デバッグや原因分析によって欠陥を見つけ、修正する必要がある。したがって、「テスト」と「デバッグ」は連続して使われることが多いが、同じ活動ではない。テストは失敗を検出する活動であり、デバッグは原因を突き止めて直す活動である。

\begin{pointbox}{重要}
テストは、欠陥の存在を示すことはできるが、欠陥が存在しないことを一般には証明できない。現実のソフトウェアでは、完全な網羅テストが不可能だからである。
\end{pointbox}

\begin{todobox}{欠陥・障害・テスト・デバッグの関係図を入れる}{白背景、モノクロ、A4 本文幅に収まる横長図を作成する。左から「欠陥」「特定の入力・状態」「障害として観測」「テストで検出」「デバッグで原因特定」「修正」の箱を並べる。欠陥が存在しても入力条件がそろわなければ障害が出ないことを点線で示す。日本語ラベル、講義ノート風。}
\end{todobox}

\subsection{主要論点}
\subsubsection{テストケース作成}
テストケース作成とは、特定の目的に役立つテストスイートを作る活動である。目的は、十分性の確認、正確性の確認、法令や標準への適合評価などである。テストケース作成は手間が大きいため、自動生成、モデルベース生成、仕様からの導出、コード構造からの導出などの技法やツールが用いられる。

\subsubsection{テスト選択と十分性基準}
テスト選択基準は、どのテストケースを選ぶかを決めるための基準である。テスト十分性基準は、どの程度まで実施すれば十分とみなすかを決める基準である。たとえば、全分岐を少なくとも一度通す、すべての機能要求を少なくとも一度検証する、主要なリスクをすべて扱う、などが考えられる。

\subsubsection{優先順位付けと最小化}
テストケース優先順位付けは、実行順序を決める活動である。重要度、リスク、過去の欠陥検出率、コードカバレッジ、類似度などを使い、価値の高いテストから先に実行する。テストケース最小化は、冗長なテストケースを削り、テストスイートを小さくする活動である。目的は、効果をできるだけ保ちながら費用と時間を下げることである。

\subsubsection{テスト目的}
テストは、目的によって選ぶ技法も評価方法も変わる。機能が仕様どおりかを確認するテスト、性能を測るテスト、セキュリティ上の弱点を探すテスト、利用者が使いやすいかを確認するテストでは、必要なテストケースも期待結果も異なる。目的を曖昧にしたままテストすると、「たくさん実行したが何を確認したかわからない」という状態になりやすい。

\subsubsection{評価と認証}
評価や認証のためのテストでは、要求、法令、標準、契約に対する根拠を示す必要がある。テストケースがどの基準から導かれたか、どの構成のシステムで実行されたか、同じ手順で再現できるかが重要である。認証が必要な領域では、テスト結果だけでなく、プロセスや文書の管理も重視される。

\subsubsection{品質保証と品質改善のためのテスト}
テストは、品質保証と品質改善にも使われる。品質保証では、要求された品質水準を満たしていることを示す根拠を作る。品質改善では、失敗、欠陥、傾向、弱点を把握し、次の設計、実装、プロセス改善につなげる。

\subsubsection{オラクル問題}
オラクル問題とは、実行結果が正しいかどうかを判断することが難しい問題である。期待結果が明確な四則演算なら判断は容易である。しかし、推薦システム、画像認識、複雑な最適化、機械学習モデル、長時間運用システムでは、何を正しい結果とみなすかが難しい。オラクルを自動化することも、しばしば高コストである。

\subsubsection{理論的・実践的制約}
テストには限界がある。代表的な限界は、完全な網羅ができないこと、成功したテストが欠陥の不存在を証明しないこと、運用環境を完全には再現できないこと、期待結果を決めるのが難しいこと、テストケースを増やすほど費用が増えることである。

\subsubsection{実行不能経路の問題}
実行不能経路とは、制御フロー上は存在するように見えるが、どの入力でも実際には通れない経路である。構造ベーステストでは、経路を通すテストを作ろうとしても、実行不能経路があると達成できない。実行不能経路の特定は、テスト時間の削減やセキュリティ上の弱点分析にも関係する。

\subsubsection{テスト容易性}
テスト容易性には二つの意味がある。第一に、対象がテストしやすい構造になっているかである。第二に、欠陥がある場合にテストスイートが障害を露出させやすいかである。観測しにくい内部状態、外部環境に強く依存する処理、再現性の低い並行処理は、テスト容易性を下げる。

\subsubsection{テスト実行と自動化}
テスト自動化は、入力生成、実行、結果比較、ログ収集、レポート作成、回帰テストなどを支援する。自動化は有効であるが、すべてを自動化すればよいわけではない。期待結果が曖昧なテスト、探索的な確認、人間の評価が必要な使いやすさの確認では、人手の役割が残る。

\subsubsection{拡張性}
拡張性は、負荷、取引数、データ量、利用者数、分散環境の複雑さが増えたときに、ソフトウェアが対応できる能力である。拡張性のテストでは、単に機能が動くかではなく、規模が増えたときの処理能力、安定性、資源使用量を確認する。

\subsubsection{テスト有効性}
テスト有効性は、テストが目的にどれだけ役立ったかを表す。欠陥検出率、最初の障害を見つけるまでのテスト数、カバレッジ、信頼性向上、残存リスク低減などで評価される。目的が違えば、有効性の測り方も変わる。

\subsubsection{制御可能性・再現性・一般化可能性}
制御可能性は、テスト条件をどの程度制御できるかである。再現性は、別の人が同じ手順で同じ結果を得られるかである。一般化可能性は、あるテスト結果が他の環境、利用者、設定にも当てはまるかである。研究でも実務でも、テスト結果を信頼するにはこの三つが重要である。

\subsubsection{オフラインテストとオンラインテスト}
オフラインテストは、外部環境との直接的な相互作用を切り離してテストする。オンラインテストは、実際の利用環境や本番に近い環境と相互作用しながらテストする。オフラインは制御しやすく、オンラインは現実に近い。どちらも期待結果を用いて評価する。

\subsection{他の活動との関係}
ソフトウェアテストは、品質管理、形式的検証、デバッグ、プログラム構築、セキュリティ、見積り、法的責任と関係する。しかし、それらと同じではない。レビューや静的解析は実行しない確認であり、テストを補完する。形式的検証は数学的に性質を示す活動であり、テストとは根拠の作り方が異なる。デバッグは、テストで観測された障害の原因を探して直す活動である。

\section{テストレベル}
テストレベルは、何を対象にテストするか、どの目的でテストするかによって分けられる。大きなシステムでは、単体、結合、システム、受入という四つの段階で考えると理解しやすい。ただし、これらは開発プロセスを固定するものではなく、どれが常に最重要というものでもない。

\subsection{テスト対象}
\subsubsection{単体テスト}
単体テストは、個別にテストできる要素を単独で確認する活動である。対象は、関数、クラス、モジュール、コンポーネント、サブシステムなどである。多くの場合、コードを書いた開発者が実施するが、必ずしもそれに限られない。単体テストの目的は、欠陥を早く見つけ、後段の結合やシステムテストに欠陥を持ち込まないことである。

\subsubsection{結合テスト}
結合テストは、複数の要素の相互作用を確認する活動である。モジュール間、コンポーネント間、外部システムとのインタフェース、ハードウェアや運用環境とのやり取りが対象になる。上位から下位へ進めるトップダウン、下位から上位へ進めるボトムアップ、混合方式、一度に結合するビッグバン方式などがある。

\subsubsection{システムテスト}
システムテストは、システム全体の振る舞いを確認する活動である。単体テストと結合テストで多くの欠陥は見つかっているはずだが、システム全体で初めて現れる問題もある。機能だけでなく、セキュリティ、プライバシー、速度、正確性、信頼性などの非機能要求の確認にも適している。

\subsubsection{受入テスト}
受入テストは、システムが要求とエンドユーザの期待を満たし、利用または配備に進めるかを確認する活動である。利用者や顧客とともに実施されることが多い。運用受入、使いやすさ確認、業務シナリオ確認などを含む。受け入れテスト駆動開発では、機能を実装する前に受入条件を定義する。

\begin{todobox}{テストレベルの四層図を入れる}{白背景、モノクロ、A4 本文幅に収まる縦長の階層図を作成する。下から「単体テスト」「結合テスト」「システムテスト」「受入テスト」の四層を積み上げる。各層に「対象」「主な目的」「典型的な実施者」を短く記載する。右側に「対象が大きくなる」「利用者視点に近づく」という矢印を置く。日本語ラベル、細線、講義ノート風。}
\end{todobox}

\subsection{テスト目的}
テスト目的は、何を確認したいかを表す。目的を明確にすると、テストケースの選び方、測り方、終了判断が決めやすくなる。

\subsubsection{適合性テスト}
適合性テストは、SUT が標準、規則、仕様、要求、設計、プロセス、慣行に合っているかを確認する。

\subsubsection{法令・規制順守テスト}
法令・規制順守テストは、SUT が法律や規制に従っていることを示す。外部の規制機関や認証制度によって求められることが多い。

\subsubsection{導入テスト}
導入テストは、SUT を対象環境へ導入した後に、ハードウェア構成、運用制約、インストール手順を含めて確認する活動である。

\subsubsection{アルファテストとベータテスト}
アルファテストは、リリース前に少数の選ばれた利用者に試してもらう活動である。ベータテストは、より広い代表的利用者に試してもらう活動である。いずれも、利用者から問題報告を受けるために行う。

\subsubsection{回帰テスト}
回帰テストは、変更後に意図しない悪影響が出ていないことを確認する再テストである。アジャイル、DevOps、テスト駆動開発、継続的開発では特に重要である。すべてのテストを毎回実行するのは高コストなので、選択、最小化、優先順位付けがよく使われる。

\subsubsection{優先順位付けテスト}
優先順位付けテストは、テストケースの実行順を工夫し、早い段階で重要な障害を見つけることを狙う。類似していないテストケースから先に実行する類似度ベースの方法などがある。

\subsubsection{非機能テスト}
非機能テストは、性能、使いやすさ、信頼性などの非機能面を確認する。代表的な種類を表に示す。

\begin{longtable}{P{0.24\linewidth}P{0.66\linewidth}}
\toprule
\textbf{種類} & \textbf{確認すること} \\
\midrule
性能テスト & 応答時間、処理能力、容量など、性能要求を満たすかを確認する。\\
負荷テスト & 多数の要求や高負荷時に、デッドロック、競合、メモリ漏れ、安定性問題が出ないかを確認する。\\
ストレステスト & 想定を超える負荷を与え、限界や破綻の仕方を確認する。\\
容量テスト & 内部保存容量、データ交換量、情報量に対する限界を確認する。\\
フェイルオーバーテスト & 障害や高負荷時にも、予備資源や代替経路でサービスを継続できるかを確認する。\\
信頼性テスト & 運用プロファイルなどを用いて、将来の信頼性を評価する。\\
互換性テスト & 異なるハードウェア、ソフトウェア、バージョン、リリースと協調できるかを確認する。\\
拡張性テスト & 利用者数、取引数、データ量が増えた場合に対応できるかを確認する。\\
弾力性テスト & クラウドなどで、計算資源、記憶資源、保存資源を増減しても能力を保てるかを確認する。\\
インフラテスト & インフラ構成要素を確認し、停止時間を減らし性能を改善する。\\
バックツーバックテスト & 複数の実装や版に同じ入力を与え、出力差を分析する。\\
回復テスト & クラッシュや災害後に再起動・復旧できるかを確認する。\\
\bottomrule
\end{longtable}

\subsubsection{セキュリティテスト}
セキュリティテストは、外部攻撃から SUT とデータを守れるかを確認する。主に機密性、完全性、可用性を扱う。誤用や悪用に対してシステムがどう振る舞うかを確認する負のテストも重要である。

\subsubsection{プライバシーテスト}
プライバシーテストは、利用者の個人データを安全に扱い、情報共有方針やプライバシー方針を守れているかを確認する。

\subsubsection{インタフェースおよび API テスト}
インタフェーステストは、コンポーネント間でデータや制御情報が正しく交換されるかを確認する。API テストでは、エンドユーザのアプリケーションが API を呼び出す状況を想定し、パラメータ、環境条件、内部データを組み合わせて確認する。

\subsubsection{構成テスト}
構成テストは、複数の利用者、環境、設定に対応するソフトウェアが、指定された構成で正しく動くかを確認する。

\subsubsection{使いやすさと人間コンピュータ相互作用テスト}
使いやすさテストは、エンドユーザがソフトウェアを学び、使い、誤りから回復できるかを確認する。機能、文書、画面、操作手順、エラー回復が対象になる。

\begin{todobox}{テスト目的の地図を入れる}{白背景、モノクロ、A4 本文幅のマインドマップを作成する。中央に「テスト目的」を置き、周囲に「適合性」「法令・規制順守」「導入」「アルファ・ベータ」「回帰」「非機能」「セキュリティ」「プライバシー」「API」「構成」「使いやすさ」を配置する。各枝に一行説明を付ける。日本語、細線、講義ノート風。}
\end{todobox}

\section{テスト技法}
テスト技法は、テストケースをどの考え方で選ぶかを定める。代表的には、仕様から選ぶ仕様ベース技法、内部構造から選ぶ構造ベース技法、経験から選ぶ経験ベース技法がある。さらに、欠陥モデル、利用モデル、アプリケーションの性質、他領域から得た知識に基づく技法もある。

\begin{todobox}{主要テスト技法の比較図を入れる}{白背景、モノクロ、A4 本文幅の比較表を作成する。列は「仕様ベース」「構造ベース」「経験ベース」「欠陥ベース」「利用ベース」。行は「主な情報源」「得意な問題」「弱点」「代表例」。各セルは短い日本語。講義ノート風、細線、印刷で読みやすい。}
\end{todobox}

\subsection{仕様ベース技法}
仕様ベース技法は、SUT の入力と出力、要求、仕様、モデルに基づいてテストケースを選ぶ。内部コードを知らなくても使えるため、ブラックボックス技法とも呼ばれる。

\subsubsection{同値分割}
同値分割は、入力領域を同じように扱えると考えられる部分集合に分け、各集合から代表値を選ぶ技法である。たとえば、年齢が 0 から 120 まで有効なら、有効範囲、下限未満、上限超過などに分ける。目的は、無数の入力を少数の代表例に縮約することである。

\subsubsection{境界値分析}
境界値分析は、入力範囲の境界付近に欠陥が集中しやすいという経験に基づき、下限、上限、その直前直後を選ぶ技法である。範囲外の値を使って堅牢性を見る場合もある。

\begin{todobox}{同値分割と境界値分析の図を入れる}{白背景、モノクロ、A4 本文幅の横長図を作成する。数直線を描き、中央に「有効範囲 0 から 120」、左右に「無効範囲」を置く。同値クラスを帯で示し、境界値として「-1」「0」「1」「119」「120」「121」を丸で示す。日本語ラベル、講義ノート風。}
\end{todobox}

\subsubsection{構文テスト}
構文テストは、形式言語で書かれた仕様や入力文法に基づいてテストケースを導く技法である。仕様が形式的であれば、機能テストケースの自動生成や期待結果の判定に使いやすい。

\subsubsection{組合せテスト技法}
組合せテストは、入力値や条件の組合せを体系的に選ぶ技法である。すべての組合せを試すと数が爆発するため、ペアワイズテストのように、任意の二つのパラメータの組合せはすべて含めるという基準を使うことが多い。

\subsubsection{決定表}
決定表は、条件と動作の組合せを表にして、論理関係を明示する技法である。複数条件によって結果が変わる業務規則に向いている。各条件組合せからテストケースを導けるため、抜け漏れの発見にも役立つ。

\subsubsection{原因結果グラフ}
原因結果グラフは、入力条件である原因と、出力や動作である結果の関係を論理グラフとして表す技法である。仕様分析、入力条件の整理、結果条件の特定に役立つ。

\subsubsection{状態遷移テスト}
状態遷移テストは、SUT を有限状態機械として表し、状態と遷移をカバーするようにテストケースを作る技法である。画面遷移、通信プロトコル、注文状態、予約状態など、状態によって振る舞いが変わる対象に向いている。

\begin{todobox}{状態遷移テストの図を入れる}{白背景、モノクロ、A4 本文幅に収まる状態遷移図を作成する。例として「注文受付」「支払待ち」「出荷準備」「出荷済み」「キャンセル」の状態を丸で示す。矢印に「支払い完了」「出荷」「キャンセル要求」などのイベントを付ける。右側に「状態カバレッジ」「遷移カバレッジ」の短い説明を置く。日本語ラベル、講義ノート風。}
\end{todobox}

\subsubsection{シナリオベーステスト}
シナリオベーステストは、利用者やシステムがたどる一連の行動をモデル化し、その流れに沿ってテストケースを作る技法である。通常の流れだけでなく、代替流、例外流も扱う。業務プロセステストやワークフローモデルを使うテストもこの考え方に含まれる。

\subsubsection{ランダムテスト}
ランダムテストは、入力領域からランダムにテストケースを選ぶ技法である。単純な自動化に向いているが、入力領域がわかっている必要がある。不正入力や予期しないデータを大量に与えるファズテストは、セキュリティ分野でよく用いられる。

\subsubsection{証拠に基づく技法}
証拠に基づく技法は、厳密な調査手順で得た知見をテストに反映する考え方である。問いを立て、利用できる証拠を集め、批判的に分析する。体系的マッピング研究や体系的レビューは、既存研究や実務知見を整理する手段である。

\subsubsection{例外強制テスト}
例外強制テストは、データ例外、操作例外、オーバーフロー、保護例外、アンダーフローなど、あらかじめ決めた例外やエラーを SUT が扱えるかを確認する。多くの場合、負のシナリオを使い、エラーメッセージや回復動作を確認する。

\subsection{構造ベーステスト技法}
構造ベース技法は、コードや設計の内部構造を情報源としてテストケースを選ぶ。ホワイトボックス技法とも呼ばれる。文カバレッジ、分岐カバレッジ、経路カバレッジ、循環的複雑度などが関係する。

\subsubsection{制御フローテスト}
制御フローテストは、文、分岐、判定、条件、条件組合せ、経路を対象に、どの部分を実行したかを確認する。最も強い基準の一つはすべての入口から出口までの経路を通す経路テストであるが、ループがあると現実的でないため、文カバレッジ、分岐カバレッジ、条件判定カバレッジなどの基準が使われる。

\subsubsection{データフローテスト}
データフローテストは、変数の定義、使用、未定義化に注目する。ある変数が定義された後、どの経路でどの使用に届くかを確認する。定義と使用の組合せを扱うため、データに関する隠れた欠陥を見つけやすい。

\subsubsection{構造ベース技法の参照モデル}
構造ベース技法では、制御フローグラフがよく使われる。ノードは文や文列を表し、辺は制御の移動を表す。この図に基づいて、どの文、分岐、経路を通したかを可視化できる。

\begin{todobox}{制御フローとデータフローの図を入れる}{白背景、モノクロ、A4 本文幅の二分割図を作成する。左は制御フローグラフで、開始、条件分岐、処理、終了をノードと矢印で示す。右は同じ流れに「x を定義」「x を使用」の注記を付け、定義から使用までのデータフローを太線で示す。日本語ラベル、講義ノート風。}
\end{todobox}

\subsection{経験ベース技法}
経験ベース技法は、仕様やコードだけではなく、テスト担当者の知識、直感、過去の障害、ドメイン理解を活用する。

\subsubsection{エラー推測}
エラー推測は、過去の欠陥履歴や経験に基づき、「ここに欠陥がありそうだ」と考えてテストケースを設計する技法である。形式的な技法では見つけにくい問題を見つけることがある。

\subsubsection{探索的テスト}
探索的テストは、学習、テスト設計、テスト実行を同時に進める活動である。事前にすべてのテストケースを決めるのではなく、観測した振る舞い、リスク、失敗の兆候に基づいて次のテストを考える。要求が詳細でないアジャイル開発などでよく使われる。

\subsubsection{その他の経験ベース技法}
アドホックテストは、技能、直感、経験に基づいてテストする。モンキーテストはランダムな操作で停止や異常を狙う。ペアテストは二人でテストを行い、一人が操作し、一人が観察する。ゲーミフィケーションは、テスト作業をゲーム要素に変換して参加を促す。クイックテストは小さなテストスイートで重大問題を早く見つける。スモークテストは、ビルド後に基本機能が動くかを確認し、本格的なテストに進める状態かを判断する。

\subsection{欠陥ベース技法とミューテーション技法}
欠陥ベーステストは、あり得る欠陥の種類を想定し、それを露出させるテストケースを作る。欠陥分類や直交欠陥分類を使うと、欠陥の意味情報を整理し、根本原因分析を短縮できる。

ミューテーションテストは、SUT をわずかに変更した変異体を多数作り、テストスイートがそれらの違いを検出できるかを評価する技法である。テストが変異体を見破れば、その変異体は殺されたという。変異体を多く殺せるテストスイートほど、欠陥検出能力が高いと考える。近年は、機械学習システムのテストで、入力を変化させて出力関係を見るメタモルフィックテストにも関連する。

\begin{todobox}{ミューテーションテストの図を入れる}{白背景、モノクロ、A4 本文幅の流れ図を作成する。左に「元のプログラム」、中央に「小さな変更を加えた変異体 A/B/C」、下に「同じテストスイートを実行」、右に「差を検出できた変異体は殺された」「検出できない変異体は生存」と示す。最後に「ミューテーションスコア」を記す。日本語ラベル、講義ノート風。}
\end{todobox}

\subsection{利用ベース技法}
利用ベース技法は、実際の運用や利用のされ方をモデル化し、その頻度や順序に沿ってテストケースを選ぶ。運用環境に近い条件で実行するため、受入テストや信頼性評価に向いている。

\subsubsection{運用プロファイル}
運用プロファイルは、入力や利用シナリオが実運用でどの頻度で起こるかを表す。運用プロファイルに基づくテストでは、観測した結果から将来の信頼性を推定する。難しさは、現実的なプロファイルを作ることにある。

\subsubsection{利用者観察ヒューリスティック}
利用者観察ヒューリスティックは、利用者がシステムやインタフェースをどれだけうまく使えるかを、統制された条件で観察する。認知的ウォークスルー、発話思考法、現場観察、質問票、面接などが使われる。

\subsection{アプリケーションの性質に基づく技法}
一般的な技法は多くのソフトウェアに使えるが、対象の性質によって追加の技法が必要になる。オブジェクト指向ソフトウェア、コンポーネントベースソフトウェア、ウェブアプリケーション、並行プログラム、通信プロトコル、リアルタイムシステム、安全重要システム、サービス指向ソフトウェア、オープンソース、組込み、クラウド、ブロックチェーン、ビッグデータ、AI・機械学習、モバイル、セキュリティ・プライバシー重視ソフトウェアなどでは、対象特有のリスクを考慮する。

\subsection{技法の選択と組合せ}
実務では、一つの技法だけで十分な品質根拠を得ることは少ない。仕様ベースと構造ベースを組み合わせれば、外部振る舞いと内部網羅を両方確認できる。リスクが高い部分には手厚い技法を使い、低リスク部分には軽量な技法を使うなど、目的、費用、リスク、組織制約に応じて選ぶ。

\subsubsection{機能的技法と構造的技法の組合せ}
シナリオベーステストは機能的観点を強く持つ。構造ベーステストは内部構造の観点を強く持つ。両者は対立ではなく補完関係である。仕様に書かれた動作を確認しつつ、実装上の分岐やデータ流れも確認することで、見落としを減らせる。

\subsubsection{決定的テストとランダムテスト}
決定的テストは、基準に従って選ばれたテストケースを実行する。ランダムテストは、確率分布に従って入力を選ぶ。信頼性テストではランダム性が役立つことがあるが、境界や特定リスクを狙う場合には決定的な選択が有効である。

\subsection{派生知識に基づく技法}
派生知識に基づく技法は、他の研究領域や文脈から得られた知識をテストに取り入れる。デジタルツイン、シミュレーション、機械学習、ゲーミフィケーション、模擬ニューラルネットワークなどが例である。テスト活動を支援し、効果を高めることを目的とする。

\section{テスト関連測度}
測度は、テスト目標が達成されたかを判断するために必要である。テスト技法は道具であり、測度はその道具が目的に役立っているかを判断する根拠である。測度は、テスト計画の最適化、進捗管理、品質分析、完了判断に使われる。

\begin{todobox}{テスト関連測度の分類図を入れる}{白背景、モノクロ、A4 本文幅に収まる二分割図を作成する。左に「SUT の評価」、右に「実施したテストの評価」を置く。左の下に「欠陥密度」「信頼性」「運用信頼性」「KPI」、右の下に「カバレッジ」「ミューテーションスコア」「欠陥注入」「相対有効性」を配置する。日本語ラベル、講義ノート風。}
\end{todobox}

\subsection{SUT の評価}
SUT の評価では、テスト結果から対象システムの状態を判断する。代表的には、欠陥の種類、欠陥密度、信頼性、運用時の障害傾向、主要業績指標が使われる。

\subsubsection{テスト計画と設計に役立つ SUT 測定}
テスト計画には、配備頻度、リードタイム、平均回復時間、変更失敗率などの測度が使われることがある。シフトレフトや DevOps では、これらの測度がテストの計画、管理、結果評価と強く結び付く。

\subsubsection{欠陥の種類・分類・統計}
欠陥分類は、どの種類の欠陥が多いか、どの品質属性に関係するかを整理する。利用性欠陥、セキュリティ脆弱性、プライバシー攻撃、サイバーリスクなど、対象に応じて分類体系は異なる。過去の発生頻度は、どこにテスト資源を集中するかを決める根拠になる。

\subsubsection{欠陥密度}
欠陥密度は、見つかった欠陥数を SUT の規模で割った値である。規模にはコード行数、機能規模、モジュール数などが使われる。欠陥密度だけで品質を完全には表せないが、比較や傾向把握には役立つ。

\subsubsection{寿命試験と信頼性評価}
信頼性評価は、テストをいつ止めるか、次のリリースに進めるかを判断するために使われる。クラウドやフォグ環境では、高い信頼性と可用性を維持することが重要であり、テスト活動もその評価に関わる。

\subsubsection{信頼性成長モデル}
信頼性成長モデルは、観測された障害と修正の履歴から、信頼性がどのように向上するかを予測する。多くのモデルでは、障害の原因となる欠陥を修正すれば信頼性が増すと仮定する。障害数に注目するモデルと、障害間隔時間に注目するモデルがある。

\subsection{実施したテストの評価}
テストそのものの評価では、テストスイートがどれだけ強いか、どの技法が有効かを比較する。カバレッジ、ミューテーション分析、欠陥注入などが用いられる。

\subsubsection{欠陥注入}
欠陥注入は、人工的な欠陥を SUT に入れ、テストスイートがそれらをどれだけ検出できるかを調べる。検出率からテスト有効性や残存欠陥数を推定しようとする。ただし、人工欠陥が本物の欠陥を代表しているか、挿入した欠陥を取り残さないかには注意が必要である。

\subsubsection{ミューテーションスコア}
ミューテーションスコアは、殺された変異体数を生成した変異体数で割った値である。値が高いほど、テストスイートが小さな変化に敏感であることを示す。ただし、変異体の生成数が多くなるため、実行コストが問題になる。

\subsubsection{技法間の比較と相対有効性}
相対有効性は、異なるテスト技法を同じ性質に対して比較する考え方である。最初の障害を見つけるまでのテスト数、欠陥検出率、信頼性改善量などが比較軸になる。

\section{テストプロセス}
テスト技法や測度は、定義されたテストプロセスに組み込まれて初めて実務で機能する。テストプロセスは、組織レベル、管理レベル、動的テストレベルの多層構造で考えられる。組織レベルでは方針や戦略を作る。管理レベルでは計画、監視、制御、完了を扱う。動的テストレベルでは設計、実装、環境構築、実行、インシデント報告を扱う。

\begin{todobox}{テストプロセスの多層図を入れる}{白背景、モノクロ、A4 本文幅の三層図を作成する。上から「組織テストプロセス」「テスト管理プロセス」「動的テストプロセス」を配置する。各層に「方針・戦略」「計画・監視・制御・完了」「設計・実装・環境・実行・報告」を記す。左に「目的・リスク・資源」、右に「根拠・報告・改善」の矢印を置く。日本語ラベル、講義ノート風。}
\end{todobox}

\subsection{実務上の考慮}
\subsubsection{態度とエゴレスプログラミング}
テストで重要なのは、障害発見を個人攻撃と見なさない態度である。欠陥が見つかることは、品質を改善する機会である。管理者は、失敗の発見と修正を歓迎する文化を作る必要がある。アジャイルやシフトレフトでは、テスターと開発者の協調が成功に不可欠である。

\subsubsection{テスト指針と組織プロセス}
組織は、テスト方針とテスト戦略を定義する。テスト方針は、テストの目的、目標、全体範囲を示す。テスト戦略は、どのようにテストするかの指針である。リスクベーステストでは、製品リスクをもとにテストの焦点と優先度を決める。

\subsubsection{テスト管理と動的テストプロセス}
テスト管理には、計画、監視、制御、完了が含まれる。動的テストプロセスには、テスト設計と実装、テスト環境の準備と保守、テスト実行、インシデント報告が含まれる。これらは、人、ツール、方針、測度を組み合わせてライフサイクルに統合される。

\subsubsection{テスト文書化}
テスト文書は、組織テスト文書、テスト管理文書、動的テスト文書に分けられる。組織テスト文書には方針や戦略が含まれる。テスト管理文書にはテスト計画、状況報告、完了報告が含まれる。動的テスト文書にはテスト設計仕様、テストケース仕様、手順仕様、データ要求、環境要求、実行ログ、インシデント報告などが含まれる。

\subsubsection{テストチーム}
テストチームの構成は、費用、納期、組織成熟度、アプリケーションの重要度によって変わる。開発者自身がテストする場合も、独立したテスト担当者が行う場合もある。シフトレフトでは、開発者、テスター、顧客が早い段階から協力し、役割の境界が薄くなることがある。

\subsubsection{テストプロセス測度}
テストプロセスでは、指定済み、実行済み、合格、不合格のテストケース数、未解決欠陥、解決済み欠陥、欠陥検出率、残存リスク、テスト進捗などを測る。これらは、プロセスリスクの管理、根本原因分析、改善に使われる。

\subsubsection{テスト監視と制御}
テスト監視と制御は、テスト活動のデータを収集し、計画どおりに進んでいるかを確認する。必要に応じて、全体計画を修正する。監視は実行と並行して行われ、関係者へ状況報告を提供する。

\subsubsection{テスト完了}
テスト完了では、どの程度テストすれば十分か、テスト段階を終えてよいかを判断する。カバレッジ、欠陥密度、信頼性推定は支援情報になるが、それだけでは十分でない。残存障害のリスクと、さらにテストを続ける費用を比較して判断する必要がある。

\subsubsection{テスト再利用性}
テスト再利用性は、テストケース、結果、環境、知識を整理し、次のテストや回帰テストで使えるようにすることである。テスト開発が高コストで複雑な場合、再利用性は特に重要である。要求や設計の変更がテストにも反映されるよう、知識リポジトリや追跡が必要である。

\subsection{テストサブプロセスと活動}
\subsubsection{テスト計画プロセス}
テスト計画では、担当者の調整、テスト目的、完了基準、設備、文書、リスク計画を定義する。組織方針、テスト段階、テスト種類、テスト環境、実行、監視、完了、報告を整える。

\subsubsection{テスト設計と実装}
テスト設計と実装では、テストレベルと選んだ技法に基づいてテストケースを作る。ツールを使う場合、ソースコード、仕様、データ構造、テスト基準などを入力し、テストスイートを生成する。場合によっては、期待結果もオラクル機能で生成できる。

\subsubsection{テスト環境の構築と保守}
テスト環境は、テストを実行するための設備、ハードウェア、ソフトウェア、ファームウェア、手順を含む。シミュレーション環境、本番に近い環境、監視・ログ機能などを整え、他のソフトウェア工学ツールと互換性を持たせる必要がある。

\subsubsection{統制された実験とテスト実行}
テスト実行は、科学的な統制実験の原則を意識すべきである。別の人が結果を再現できるほど、手順、環境、SUT の版、入力、観測結果を明確に記録する。受入テストでは、A/B テストのように利用者の選好を統計的に評価することもある。

\subsubsection{テストインシデント報告}
テストインシデント報告は、いつ、誰が、どの構成でテストを行い、どの結果が得られたかを記録する。予期しない結果は問題報告システムに記録され、後のデバッグや変更管理の入力になる。すべての異常が欠陥とは限らないため、分析により原因を特定する必要がある。

\subsection{人員配置}
人員配置では、役割、活動、責任、採用、教育を決める。スクラムマスター、テストリード、品質保証担当、テスト分析者、テスト設計者、セキュリティテスト担当、性能テスト担当、環境担当、自動化担当などの役割があり得る。必要な技能が不足すれば、欠陥発見能力、変更対応、納期、保守費用に影響する。

\section{開発プロセスと適用領域におけるテスト}
テストは、どの開発プロセスでも基本活動である。ただし、実施の位置づけや用語は、ライフサイクルや適用領域によって変わる。

\subsection{ソフトウェア開発プロセス内のテスト}
\subsubsection{従来型プロセスにおけるテスト}
逐次型、V 字型、スパイラル、反復型などの従来型プロセスでは、テストが一つの活動として明確に位置づけられることが多い。後半で集中的にテストする場合、利用者ニーズとのずれや評価上の問題が遅く見つかるリスクがある。品質を評価し制御するために、TMMi、SPI、SPICE、CMMI、統一プロセスなどの成熟度モデルや改善枠組みが用いられてきた。

\subsubsection{シフトレフトに沿ったテスト}
シフトレフトは、開発の早い段階でテストを採用し、欠陥を早く検出して除去する考え方である。アジャイル、DevOps、TDD はこの動きに含まれる。内部コード品質では、単体・結合テストや構造ベース技法が重要である。業務ニーズでは、システム・受入テスト、利用ベース技法、シナリオベース技法が重要である。知覚品質では、アルファ、ベータ、導入、使いやすさ、セキュリティ、プライバシーが重視される。品質保証では、性能、セキュリティ、プライバシー、適合性、順守が対象になる。

アジャイル開発では、顧客やチームメンバーを含む関係者がテスト活動に関与し、回帰欠陥のリスクや変化する要求への対応が重要になる。TDD では、コードを書く前にテストケースを作り、要求の理解を明確化する。DevOps や継続的インテグレーションでは、回帰テストやスモークテストを継続的に行い、配備前に基本的なテスト可能性を確認する。

\begin{todobox}{従来型とシフトレフトの比較図を入れる}{白背景、モノクロ、A4 本文幅の比較図を作成する。上段に従来型の流れ「要求、設計、実装、後半のテスト」を横矢印で示す。下段にシフトレフトの流れ「要求からテスト設計、単体、結合、回帰、自動化、監視」が早期から並走する様子を示す。右側に「欠陥を早く見つけるほど修正コストが下がる」と注記する。日本語ラベル、講義ノート風。}
\end{todobox}

\subsection{適用領域におけるテスト}
適用領域ごとに、重視すべきリスクや標準が異なる。自動車領域では、ソフトウェア、ハードウェア、信号、統合、安全、ストレス、セキュリティが関係する。IoT では、機器管理、異種性、通信、データ、プライバシー、サイバーセキュリティが重要である。法務領域では、高度に機微なデータ、用語、専門家の関与、整合性、順守が重要である。

モバイル領域では、画面解像度、GPS、OS、端末メーカー、ネイティブアプリとウェブアプリの違いを考慮する。航空電子領域では、独立または疎結合の構成要素、市販品、機能・非機能、運用、通信、安全、セキュリティが関係する。医療領域では、安全で信頼できるデータ交換、性能、プライバシー、安全、規制順守が重要である。

組込み領域では、ソフトウェアとハードウェアが密接に結合する。GUI では、画面要素、見た目、配置、誤字、操作感が対象になる。ゲームでは、プレイテスト、機能、セキュリティ、適合性、性能が重要である。リアルタイム領域では、時間制約と決定的振る舞いが重要である。サービス指向、金融、クラウド、AI なども、それぞれ特有の観点を持つ。

\section{新興技術のテストと新興技術によるテスト}
新興技術は、テスト対象として新しい課題を生む一方で、テスト活動を支援する手段にもなる。この章では、「新興技術をテストする」と「新興技術を使ってテストする」を分けて理解する。

\begin{todobox}{Testing of と Testing through の対比図を入れる}{白背景、モノクロ、A4 本文幅の二分割図を作成する。左に「新興技術をテストする」として AI、ブロックチェーン、クラウド、分散アプリを配置する。右に「新興技術でテストする」として ML、ブロックチェーン、クラウド、シミュレーション、クラウドソーシングを配置する。中央に「対象か、手段か」という吹き出しを置く。日本語ラベル、講義ノート風。}
\end{todobox}

\subsection{新興技術をテストする}
AI、機械学習、深層学習を使ったシステムでは、非決定性、データ依存性、オラクル問題が大きな課題になる。欠陥は、学習データ、学習プログラム、モデル、フレームワーク、推論環境のどこにでもあり得る。オフラインでは交差検証などでモデルを評価し、配備後はオンラインテストで利用者行動との相互作用を評価する。

ブロックチェーンでは、スマートコントラクトや関連アプリケーションに対して、ストレステスト、侵入テスト、性質テストが使われる。公開型か私有型か、他アプリケーションとの接続、トランザクション性能、セキュリティなどを考慮する。

クラウドでは、クラウド上に配備されたアプリケーションやインフラの機能・非機能品質を確認する。性能、拡張性、弾力性、セキュリティ、互換性、相互運用性が重要である。分散・並行アプリケーションでは、複数 OS、ブラウザ、ハードウェア、利用者、依存関係、継続的配備がテストを複雑にする。

\subsection{新興技術を使ってテストする}
機械学習は、テストケース設計、オラクル問題、テスト評価、優先順位付け、ミューテーションテスト自動化、欠陥分類、トリアージに使われる。大量データから傾向を抽出し、保守労力を減らす可能性がある。

ブロックチェーンは、分散した関係者が信頼できる形でテスト活動を調整する手段になり得る。改ざん耐性、監査可能性、分散データ管理、要求順守の自動確認に役立つ可能性がある。クラウドは、大規模シミュレーションや弾力的資源を必要とするテスト環境として利用できる。

シミュレーションは、危険な状況、災害、復旧、特定振る舞いを評価するための重要な手段である。自動車や組込み領域では、実信号を SUT に与えながら現実を模擬するハードウェアインザループテストが使われる。クラウドソーシングテストは、多様な利用者、場所、端末、条件での検出力を得るために使われるが、社内検証の代替ではなく補完である。

\section{ソフトウェアテストツール}
テストは、多数の実行、結果の比較、大量の情報管理を伴うため、ツール支援が重要である。ツールは、単純な反復作業を減らし、テスト設計や生成を支援し、効率と有効性を高める。ただし、目的、開発方式、評価対象、実行環境に合うツールを選ぶ必要がある。一つの万能ツールではなく、複数のツールを組み合わせることが多い。

\subsection{テストツール支援と選択}
ツール選択では、何を自動化したいか、どの測度を集めたいか、どの環境で実行するか、既存プロセスとどう統合するかを確認する。ツール導入自体にも学習コスト、保守コスト、環境依存、誤判定、オラクルの限界がある。

\subsection{ツールの分類}
代表的なツール分類を表に示す。

\begin{longtable}{P{0.27\linewidth}P{0.63\linewidth}}
\toprule
\textbf{分類} & \textbf{主な役割} \\
\midrule
テストハーネス & ドライバやスタブにより、部分的な SUT を制御された環境で実行し、出力を記録する。\\
テスト生成器 & ランダム、経路ベース、モデルベースなどによりテストケースを生成する。\\
キャプチャ・リプレイツール & 以前の入力や画面操作を記録し、再実行する。\\
オラクル・比較・表明検査ツール & 実行結果が成功かどうかの判定を支援する。\\
カバレッジ分析器・計測器 & プログラム中の文、分岐、経路などがどれだけ実行されたかを測る。\\
トレーサ & 実行経路の履歴を記録する。\\
回帰テストツール & 変更後の再実行、影響範囲に基づくテスト選択を支援する。\\
信頼性評価ツール & テスト結果を分析し、信頼性関連測度を可視化する。\\
注入ベースツール & 攻撃注入や欠陥注入により、特定条件での振る舞いを確認する。\\
シミュレーションベースツール & モデルを使ってシナリオを実行し、性質の検証や予測を行う。\\
セキュリティテストツール & 脆弱性、侵入、ファズテストなどを支援する。\\
テスト管理ツール & テスト計画、実行、データ収集、欠陥管理、報告を支援する。\\
クロスブラウザテストツール & 端末やブラウザごとに UI が期待どおり動くかを確認する。\\
負荷テストツール & 性能評価のために負荷を発生させ、関連データを集める。\\
欠陥追跡ツール & 欠陥報告、状態、修正履歴を管理する。\\
モバイルテストツール & 実機やエミュレータで、モバイルアプリの UI や動作を繰り返し確認する。\\
API テストツール & API の機能、性能、信頼性、セキュリティを自動テストする。\\
ウェブアプリケーションテストツール & ウェブベース SUT の機能、性能、利用者向け問題を確認する。\\
\bottomrule
\end{longtable}

\begin{todobox}{テストツール分類図を入れる}{白背景、モノクロ、A4 本文幅の分類図を作成する。中央に「テストツール」を置き、周囲に「生成」「実行」「判定」「測定」「管理」「セキュリティ」「性能」「環境・端末」を配置する。各分類に代表例を 2 から 3 個ずつ付ける。日本語ラベル、講義ノート風。}
\end{todobox}

\section{トピックと参考文献の対応}
この表は、SWEBOK が示す「トピック対参考資料の行列」を、学習用に簡略化したものである。詳しく学ぶときは、各トピックに対応する文献を読む。

\begin{longtable}{P{0.46\linewidth}P{0.46\linewidth}}
\toprule
\textbf{トピック} & \textbf{主な参考資料} \\
\midrule
1. ソフトウェアテストの基礎 & Naik and Tripathy; Sommerville; Laporte and April \\
1.1 欠陥と障害 & Naik and Tripathy; Sommerville; Laporte and April \\
1.2 主要論点 & Naik and Tripathy; ISO/IEC/IEEE 29119; IEEE 1012; ISO/IEC 25010; ISO/IEC/IEEE 32675 \\
1.3 他活動との関係 & Software Quality KA; Software Construction KA; Software Security KA; Professional Practice KA \\
2. テストレベル & Naik and Tripathy; Sommerville \\
2.1 テスト対象 & Naik and Tripathy; Sommerville \\
2.2 テスト目的 & Naik and Tripathy; Sommerville; ISO/IEC/IEEE 29119; ISO/IEC/IEEE 24765; Nielsen \\
3. テスト技法 & Naik and Tripathy; ISO/IEC/IEEE 29119; Utting et al.; Nielsen \\
3.1 仕様ベース技法 & Naik and Tripathy; ISO/IEC/IEEE 29119; Utting et al. \\
3.2 構造ベース技法 & Naik and Tripathy; ISO/IEC/IEEE 29119 \\
3.3 経験ベース技法 & ISO/IEC/IEEE 29119; Riccio et al. \\
3.4 欠陥ベース・ミューテーション & Naik and Tripathy; Papadakis et al.; ISO/IEC/IEEE 24765 \\
3.5 利用ベース技法 & Naik and Tripathy; Sommerville; Nielsen \\
3.6 アプリケーションの性質 & Sommerville; Laporte and April; IEEE 1012 \\
3.7 技法の選択と組合せ & Laporte and April; ISO/IEC/IEEE 32675; ISO/IEC/IEEE 29119 \\
3.8 派生知識に基づく技法 & Sommerville; Laporte and April; Utting et al. \\
4. テスト関連測度 & Sommerville; Laporte and April; ISO/IEC/IEEE 32675; ISO/IEC/IEEE 29119 \\
5. テストプロセス & Sommerville; ISO/IEC/IEEE 29119; SWECOM; ISO/IEC 20246 \\
6. 開発プロセスと適用領域 & Sommerville; Laporte and April; ISO/IEC/IEEE 29119 \\
7. 新興技術のテストと新興技術によるテスト & Riccio et al.; Demi et al.; Bertolino et al.; Achary and Raj; Mao et al. \\
8. ソフトウェアテストツール & Naik and Tripathy; Laporte and April \\
\bottomrule
\end{longtable}

\section{参考文献}
\begin{enumerate}[leftmargin=2.2em]
\item S. Naik and P. Tripathy, \textit{Software Testing and Quality Assurance: Theory and Practice}, 1st ed., Wiley, 2008.
\item I. Sommerville, \textit{Software Engineering}, 10th ed., Addison-Wesley, 2016.
\item E. W. Dijkstra, \textit{Notes on Structured Programming}, Technological University, Eindhoven, 1970.
\item ISO/IEC/IEEE 29119, \textit{System and software engineering - Software testing}, 2022.
\item ISO/IEC/IEEE 24765:2017, \textit{Systems and Software Engineering - Vocabulary}, 2nd ed., 2017.
\item M. Papadakis, M. Kintis, J. Zhang, Y. Jia, Y. Le Traon, and M. Harman, ``Mutation Testing Advances: An Analysis and Survey,'' \textit{Advances in Computers}, 112, 2019.
\item M. Utting, B. Legeard, F. Bouquet, E. Fourneret, F. Peureux, and A. Vernotte, ``Recent advances in model-based testing,'' \textit{Advances in Computers}, 101, 2016.
\item IEEE Std 1012-2016, \textit{IEEE Standard for System, Software, and Hardware Verification, and Validation}, 2016.
\item ISO/IEC 25010:2011, \textit{Systems and software Quality Requirements and Evaluation - System and Software Quality Models}, 2011.
\item ISO/IEC/IEEE 32675:2022, \textit{Information technology - DevOps - Building reliable and secure systems including application build, package and deployment}, 2022.
\item Software Engineering Competency Model (SWECOM), v1.0, 2014.
\item ISO/IEC 20246:2017, \textit{Software and systems engineering - Work product reviews}, 2017.
\item V. Riccio, G. Jahangirova, A. Stocco, et al., ``Testing machine learning based systems: A systematic mapping,'' \textit{Empirical Software Engineering}, 25, 2020.
\item C. Y. Laporte and A. April, \textit{Software Quality Assurance}, IEEE Computer Society Press, 1st ed., 2018.
\item S. Demi, R. Colomo-Palacios, and M. Sánchez-Gordón, ``Software Engineering Applications Enabled by Blockchain Technology: A Systematic Mapping Study,'' \textit{Applied Sciences}, 11(7), 2021.
\item K. Mao, L. Capra, M. Harman, and Y. Jia, ``A survey of the use of crowdsourcing in software engineering,'' \textit{Journal of Systems and Software}, 126, 2017.
\item A. Bertolino, G. D. Angelis, M. Gallego, B. García, F. Gortázar, F. Lonetti, and E. Marchetti, ``A systematic review on cloud testing,'' \textit{ACM Computing Surveys}, 52(5), 2019.
\item R. Achary and P. Raj, \textit{Cloud Reliability Engineering: Technologies and Tools}, CRC Press, 2021.
\item J. Nielsen, \textit{Usability Engineering}, 1st ed., Morgan Kaufmann, 1993.
\end{enumerate}

\section{画像 TODO 一覧}
本文に配置した画像 TODO を再掲する。実際に図を作る場合は、各 TODO のプロンプトをそのまま画像生成に使うと、A4 資料に合う図を作りやすい。

\begin{enumerate}[leftmargin=2em]
\item 05 章の全体地図。
\item テストの基本定義図。
\item 欠陥・障害・テスト・デバッグの関係図。
\item テストレベルの四層図。
\item テスト目的の地図。
\item 主要テスト技法の比較図。
\item 同値分割と境界値分析の図。
\item 状態遷移テストの図。
\item 制御フローとデータフローの図。
\item ミューテーションテストの図。
\item テスト関連測度の分類図。
\item テストプロセスの多層図。
\item 従来型とシフトレフトの比較図。
\item Testing of と Testing through の対比図。
\item テストツール分類図。
\end{enumerate}

\section{章末確認問題}
\begin{enumerate}[leftmargin=2em]
\item 欠陥と障害の違いを説明せよ。なぜテストは欠陥ではなく障害を直接観測することが多いのか。
\item テストケースに入力値以外の情報が必要な理由を説明せよ。
\item テストオラクルとは何か。オラクル問題が起きやすいシステムの例を一つ挙げよ。
\item 単体テスト、結合テスト、システムテスト、受入テストを、対象と目的の観点から比較せよ。
\item 回帰テストがアジャイルや DevOps で重要になる理由を説明せよ。
\item 同値分割と境界値分析の違いを、年齢入力の例で説明せよ。
\item 仕様ベース技法と構造ベース技法はなぜ補完関係にあるのか。
\item 探索的テストが有効になる状況を説明せよ。
\item ミューテーションスコアは何を表すか。値が高いと何が言えるか。
\item テスト完了をカバレッジだけで判断してはいけない理由を説明せよ。
\item シフトレフトの考え方を、従来型プロセスと比較して説明せよ。
\item 「新興技術をテストする」と「新興技術を使ってテストする」の違いを説明せよ。
\end{enumerate}

\begin{pointbox}{解答の方向性}
1 は、原因としての欠陥と観測される結果としての障害を区別する。2 は、状態、環境、手順、期待結果が振る舞いに影響することを述べる。3 は、結果判定の仕組みであり、機械学習や推薦などで難しい。4 は対象の大きさと利用者視点への近さで整理する。5 は変更が頻繁で、既存機能の破壊を早く見つける必要があることを述べる。6 以降は、本文の「選択基準」「十分性」「カバレッジ」「リスク」「費用」「再現性」「目的に応じた技法選択」という語を使って説明できる。
\end{pointbox}

\end{document}
